\chapter{Bill of materials}

\section{Essential components}

The following components are indispensable for the system in its minimum configuration (a single phone). Prices are indicative, in Chilean pesos as of mid-2026.

\begin{longtable}{p{0.4\textwidth} p{0.32\textwidth} R{0.18\textwidth}}
\toprule
\textbf{Component} & \textbf{Suggested model} & \textbf{Price (CLP)} \\
\midrule
\endhead
Single Board Computer & Raspberry Pi 4 Model B (2\,GB) & \$85,000 \\
Power supply & Official Raspberry Pi USB-C 15\,W (5\,V/3\,A) & \$15,000 \\
Storage & MicroSD SanDisk 32\,GB Class 10 / A1 & \$10,000 \\
Thermal dissipation & Case with fan for Pi 4 & \$12,000 \\
Video cable & micro-HDMI to HDMI cable, 1.5\,m & \$5,000 \\
Telephone adapter & Grandstream HT801 (1 FXS port) & \$40,000 \\
Telephone cable & RJ11 cable, 2\,m & \$2,000 \\
Network cable & Ethernet Cat5e patch cord, 1\,m & \$3,000 \\
Analog telephone & Any landline phone with DTMF (1990s+) & --- \\
Television/monitor & Any device with HDMI input & --- \\
\midrule
\textbf{Essentials total} & & \textbf{\textasciitilde\$172,000} \\
\bottomrule
\caption{Essential bill of materials.}\\
\end{longtable}

\subsection{Per-component details}

\subsubsection{Raspberry Pi 4 (2\,GB)}

Available in Chile through official distributors such as \texttt{raspberrypi.cl} and \texttt{mcielectronics.cl}. The Pi 4B model with at least 2\,GB of RAM is required. Older models (3B+ or below) don't have enough emulation performance, and models with less RAM leave the system with no headroom.

\begin{note}
If the Pi 4 (2\,GB) is out of stock, the Pi 5 (2\,GB) works as well and performs better. The price difference is usually around 10\%. Either one meets the requirements.
\end{note}

\subsubsection{Official USB-C power supply}

\textbf{Critical.} The Pi 4 demands a steady 5\,V at 3\,A under load. Mobile phone chargers, even branded ones, often deliver unstable voltage, which causes \emph{undervoltage warnings}, spontaneous reboots and microSD corruption. The official supply is the only reliable way to avoid it.

\subsubsection{MicroSD}

32\,GB is more than enough. The important specifications are:

\begin{itemize}
  \item \textbf{Class 10}: guarantees 10\,MB/s sustained write.
  \item \textbf{A1 or A2 endurance}: optimized for applications (not just large files such as video).
\end{itemize}

Recommended brands: SanDisk Endurance, Samsung Evo Plus.

\begin{tip}
Cheap unbranded microSDs fail within months when an operating system writes constantly. For an always-on \emph{appliance}, it's worth spending a little more.
\end{tip}

\subsubsection{Grandstream HT801}

An \emph{Analog Telephone Adapter} (ATA) with one FXS \emph{(Foreign eXchange Subscriber)} port, that is, a port that supplies line voltage (\textasciitilde48\,V DC) and ringing voltage (\textasciitilde90\,V AC) to an analog telephone, and on the other side connects to the local Ethernet network and registers as a SIP client.

Relevant features:

\begin{itemize}
  \item G.711 (PCMU/PCMA), G.722, G.729, iLBC, OPUS codecs
  \item In-band DTMF, RFC 2833 and SIP INFO
  \item Web-configurable
  \item \emph{Off-hook Auto-Dial} feature: automatically dials a number when the handset is lifted
\end{itemize}

In Chile it's sold via MercadoLibre, VoIP distributors such as Inttelec or Centralinet, or by direct import from Aliexpress (with warranty risk).

\subsubsection{Analog telephone}

Any landline phone with a DTMF keypad from the last forty years works. \textbf{Rotary-dial phones don't work} because they generate pulses, not tones. Classic models that fit the role and add 1990s aesthetic:

\begin{itemize}
  \item Panasonic KX-T series
  \item General Electric with curly cord
  \item Early cordless VTech (with a large base)
  \item Any ``hotel'' phone Cetis or similar
\end{itemize}

\section{Optional components}

\begin{longtable}{p{0.4\textwidth} p{0.4\textwidth} R{0.15\textwidth}}
\toprule
\textbf{Component} & \textbf{Reason} & \textbf{Price (CLP)} \\
\midrule
\endhead
Small used CRT TV & Authentic 1990s aesthetic & \$20,000--40,000 \\
5-port Ethernet switch & If the router is far from the setup & \$10,000 \\
USB HDMI capture card & To record/stream gameplay & \$25,000 \\
External 3.5\,mm speakers & If the TV has weak audio & \$15,000 \\
Decorative labels & ``Player 1'' style design on the phone & \$3,000 \\
\bottomrule
\caption{Optional components.}\\
\end{longtable}

\section{Growth plan: multi-player}

The base system supports a single phone. To scale up to the TV mode with multiple simultaneous players:

\begin{itemize}
  \item Replace the HT801 with an \textbf{HT802} (2 FXS ports, $\sim$\$55,000) or \textbf{HT814} (4 FXS ports, $\sim$\$110,000).
  \item Acquire additional phones.
  \item Reconfigure Asterisk with \texttt{ConfBridge} and turn-queue logic.
\end{itemize}

This growth is documented in detail in chapter 9.

\section{Comparison with alternatives}

\begin{table}[h]
\centering
\begin{tabular}{lccc}
\toprule
& \textbf{Pi 4 (2\,GB)} & \textbf{Pi 5 (2\,GB)} & \textbf{Used x86 mini PC} \\
\midrule
Approx. price (CLP) & \$85,000 & \$95,000 & \$100,000 \\
Hugo emulation & Good & Very good & Excellent (native) \\
Power draw & 5--7\,W & 8--12\,W & 15--25\,W \\
Asterisk & Perfect & Perfect & Perfect \\
Physical size & Credit card & Credit card & Small \\
HDMI to TV & Yes (micro) & Yes (micro) & Yes (standard) \\
Appliance look & Excellent & Excellent & Good \\
\bottomrule
\end{tabular}
\caption{Platform comparison for Hugo Phone.}
\end{table}

For this project the \textbf{Pi 4 (2\,GB)} is selected for its combination of price, size and power efficiency.
